

	A final option is, like previous deep learning successes [193, 24], to leverage the vast quantities of data available on the internet.

	. Specifically, our interest lies in internet video data. Our reasoning here is threefold:
(1) Relevant information content. In contrast to internet-scraped text or image data, internet video can uniquely
offer information regarding the physics and dynamics of the world, and information regarding human
behaviours and actions [308]. Crucially, internet video has coverage over many behaviours and tasks
relevant to a general-purpose robot (e.g., household chores).
(2) High in quantity and diversity. There are huge quantities of video data freely available on the internet [249].
Importantly, this data is highly diverse. The largest open-source robot dataset [195] pales in comparison,
both in terms of quantity and diversity of the data (see Figure 8).
(3) Internet video is relatively untapped. The use of real and simulated robot data has been extensively explored
[36, 264, 8, 127]. Meanwhile, leveraging pretrained text and image foundation models has become increasingly common in robotics [6, 157, 37, 238]. However, the use of internet video data in robotics is in more
nascent stages.

• Internet data. Robot learning has benefited from the use of broad internet data. Image and video data have
been used to pretrain visual representations for robotics [277, 188]. Foundational VLMs and LLMs have been
used to help define reward functions for the robot learner [258, 69, 317, 133]. LLMs have been employed
as high-level planners in long-horizons tasks [6, 112]. Some works have finetuned a standard internetpretrained foundation model into a robot foundation model [37, 131]. Others have jointly pretrained on
both internet data and action-labelled robot data [217, 251]. We elaborate on how internet video data has
been used to aid robot learning in Section 4.3.


- our interest lies in internet video data
    - **Relevant information content.**
        - offer information regarding the physics and dynamics of the world and information regarding human behaviours and actions [308]
        - many behaviours and tasks relevant to a general-purpose robot
    - **High in quantity and diversity.**
    - **Internet video is relatively untapped.**


	- Internet data
	    - Image and video data have been used to pretrain visual representations for robotics [277, 188].
    - Foundational VLMs and LLMs have been used to help define **reward functions** for the robot learner [258, 69, 317, 133].
    - **LLMs** have been employed as **high-level planners in long-horizons tasks** [6, 112].
    - Some works have **finetuned a standard internet-pretrained foundation model** into a robot foundation model [37, 131].
    - Others have jointly pretrained on both internet data and action-labelled robot data [217, 251].























	\item \texttt{Simulated data}
	
	Another option is to leverage simulation
	[336, 127]. However, simulation comes with issues related to flawed simulated physics and difficulties creating
	and training on a suitable diversity of simulated environments and tasks.

	- Simulated data
    - use simulations
    - Inaccuracies in low-level physics create a **‘sim-to-real’ gap** [329]
    **—> solution here is to employ domain randomization [266]
    - challenge of Manually creating a suitable diversity of simulated environments and tasks for generalist robotics 
    **—> solve this problems using procedural environment generation [62]
    - policy lack to collect high-quality data in the simulation.
    —> use of an automated curriculum during RL training [200, 158] 
    —> collecting data using specialist policies [259]
    —> models with access to privileged simulation information [92, 276, 57]
    —> bootstrapping data collection from human demonstrations [179, 177].


The use of simulation is promising for scaling up online learning while avoiding difficulties associated with the real-world [336, 127, 190]. However, the use of simulation presents a number
of issues [25]. (1) Inaccuracies in low-level physics create a ‘sim-to-real’ gap [329] that must be overcome.
A partial solution here is to employ domain randomization [266]. (2) Manually creating a suitable diversity
of simulated environments and tasks for generalist robotics is a challenge. Recent works seek to tackle this
using procedural environment generation [62, 262] and foundation-model-assisted environment design
[292, 76, 286]. (3) We often lack a policy to collect high-quality data in the simulation. Solutions here
have included the use of an automated curriculum during RL training [200, 158] and collecting data
using: specialist policies [259], models with access to privileged simulation information [92, 276, 57], or
bootstrapping data collection from human demonstrations [179, 177].

	
	
	



	% Real-world data
	One possibility is to use humans to collect robot data [36, 129]. However, this is expensive and can be difficult in tasks requiring skill.

	Data collected via human teleoperation [36, 129] and pooled data from multiple academic labs [195] have been used to train initial robot foundation models [36, 264]. Other works have investigated 	methods for automating data collection to improve scalability [35, 5, 304]. Several companies have demonstrated evidence of infrastructure suitable for large-scale robot data-collection [251, 115]. However, the largest open-source robot dataset [195] is still significantly smaller and less diverse than internet-scale data
	(see Figure 8).
	
	- Real world data
		- **Data collected via human teleoperation** [36, 129] and pooled data from multiple academic labs [195] have been used to train initial robot foundation models [36, 264].
		- methods for **automating data collection** to improve scalability [35, 5, 304].








To train a generalist robot, large amounts of diverse data are required in order to capture the variability of real-world environments and tasks. In particular, robot interaction data, also known as robot data, is highly valuable because it contains both observations of the environment, such from RGB-D cameras or other sensors, and the corresponding robot trajectories and actions. This enables the robot to learn the relationship between observations, actions, and their resulting outcomes.

However, the currently available amount of robot data is still insufficient in both scale and diversity to train highly generalizable robotic systems. Collecting robot data, whether in real-world environments or in simulation, is expensive and time-consuming. Real-world data collection requires physical hardware, human supervision, and extensive interaction time, while simulation environments often suffer from the sim-to-real gap and still require significant engineering effort. As a result, data scarcity remains one of the central challenges in the development of general-purpose robotic systems.

To address this data bottleneck, a promising approach is to leverage large-scale internet video data. Internet videos provide a rich source of information about human behaviors, object interactions, and task execution in diverse environments. They contain fundamental information about physical dynamics, motion patterns, and semantic task structure that can be useful for robotic learning. In contrast to robot datasets, internet-scale video data is abundant, inexpensive to obtain, and covers a much broader range of scenarios and activities.

Although such videos typically do not contain direct robot actions, recent work has shown that they can still provide strong supervisory signals for representation learning, action understanding, and policy pretraining. By augmenting traditional robot datasets with large-scale video data, robotic systems can learn more robust and transferable representations, improving their ability to generalize to unseen tasks and environments \cite{lin2024data}.